% Copyright (c) 2026 Eugene Joseph M. Ragasa.
% SPDX-License-Identifier: Apache-2.0
% Author: Eugene Joseph M. Ragasa
% Formal proof completed: 2026-10-19
\documentclass[11pt]{article}
\usepackage[T1]{fontenc}
\usepackage{amsmath,amssymb,booktabs,tabularx}
\usepackage[margin=1in]{geometry}
\title{Proof Dossier: Exact Retained--Omitted Decomposition}
\author{Eugene Joseph M. Ragasa}
\date{Formal proof completed: 2026-10-19}
\begin{document}\maketitle
\section{Identity and status}
\begin{tabularx}{\textwidth}{@{}lX@{}}\toprule
Statement ID & \texttt{BG-FOURIER-THM-004} \\
Status & proved; no admissions \\
Lean declaration & \texttt{ProjectKoios.BandGauge.Fourier.}\newline\texttt{finiteReconstruction\_eq\_truncated\_add\_omitted} \\
Owning module & \texttt{ProjectKoios/BandGauge/Fourier/FiniteReconstruction.lean} \\
Axioms & \texttt{propext}, \texttt{Classical.choice}, \texttt{Quot.sound} \\
\bottomrule\end{tabularx}
\section{Statement}
The full finite reconstruction is the sum of its retained-support reconstruction and its explicit complement sum.
\section{Argument}
The finite mode inventory is partitioned into the support and its complement. Finite-sum additivity gives the equality exactly.
\section{Boundary and provenance}
Completeness refers only to the declared finite type, not an external mesh or continuum domain. Lean checked the exact statement; the dossier and proof were AI-assisted.
\end{document}
